\subsection{Balancing AI Autonomy and Alignment with Humane Values}
\label{sec:5.6.3}
Autonomy and alignment pull in opposite directions almost by definition: the more independently a system acts, the less each individual action can be checked against human judgment before it happens. Closing that gap without collapsing autonomy into constant supervision is the operational version of the question section~\ref{sec:5.6.2} left open — and it is the same shape of trade-off section~\ref{sec:5.3} covers at the level of a single judgment, fast and autonomous against slow and deliberate.

\runin{Where each alignment method leaves the latitude question}
Cooperative inverse reinforcement learning, introduced at section~\ref{sec:4.2.3}, matters here for one structural consequence. An agent that treats the human reward as unknown and to be inferred through interaction keeps its model of what it is optimizing open to revision for as long as the interaction lasts, which is a different shape of autonomy from an agent executing a reward function fixed at training time and never revisited. The cost of that openness is that what the agent ends up optimizing for is determined by whoever it interacts with, which is why section~\ref{sec:3.3} treats a system's correctability and its capacity to hold a line as one property.

Autonomy also needs boundaries on what a system is allowed to try while it is still learning. Constrained policy optimization lets a reinforcement-learning agent search for better policies while remaining inside predetermined constraints throughout the search and at the end of it \autocite{achiam2017constrained}. That is safe exploration, and the alternative is exploration followed afterward by a safety check.

\runin{What autonomy needs to stay inside antifascist bounds}
Autonomy that is not grounded in antifascist values is a capability without a governor. Building autonomy that stays inside those bounds means designing for fairness, transparency, and accountability from the start, since an audit only finds their absence afterward, and testing the design under adversarial pressure: red-teaming a system, deliberately simulating the ways it could be misused, before deployment, while the design can still absorb the answer. It also means keeping people other than a system's own developers in the design process — ethicists, users, and representatives of the communities most exposed to that system's failure modes — since the group that built a system is a poor judge, on its own, of whose freedoms it might end up infringing.

\runin{Every safeguard here is external, and autonomy is the distance from it}
The safeguards usually proposed here — third-party audit, disclosure requirements, export due diligence, cross-sector standards bodies, AI literacy — are the subject of chapters~\ref{sec:8} through \ref{sec:10}, which work out what each one costs and which of them survive an institution that would rather not comply. What belongs here is the part that is specific to autonomy, and it is one observation. Every safeguard on that list is external: it works by giving some outside body the standing and the information to intervene. Autonomy is the property that puts distance between the system's action and that body's attention, and the safeguards do not scale with it. A content-moderation system running without continuous human review is autonomous in the straightforward sense, and the same independence that makes it efficient lets it drift into over-censorship, or be steered toward favoring one ideology's speech over another's, to be discovered afterward by an audit reading the record of what already happened.

That leaves an uncomfortable division of labor and it is the one this book keeps arriving at. External mechanisms find out what a system did, after it did it. Declining to do it in the first place is available only to something inside the system, and section~\ref{sec:5.6.2} says why that inside thing cannot be assembled out of an outside check, while section~\ref{sec:3.3} says why it cannot be assembled out of correctability either.

None of that substitutes for a genuine override. A human able to intervene in a high-stakes case, and a system built to accept that intervention rather than resist it, is what keeps autonomy answerable to something outside itself. That is one of three places this book arrives at the same requirement from a different direction; chapter~\ref{sec:3} collects them and asks what it would take to build the thing they all point at — and concludes, at section~\ref{sec:3.3}, that this requirement and the floor are the same property with the sign flipped.
